\begingroup\setbox0=\vbox{\section*{Số tự nhiên}}\endgroup
\chapterhead{1}{Số tự nhiên}

\shead{Dãy số cách đều}
\begin{knowbox}\begin{itemize}[leftmargin=1.4em,itemsep=3pt]
\item \textbf{Dãy số cách đều} là dãy số mà kể từ số hạng thứ hai, mỗi số hạng đều lớn hơn (hoặc bé hơn) số hạng đứng liền trước nó cùng một số đơn vị, gọi là \textbf{khoảng cách}.
\item Ví dụ: $3,\ 7,\ 11,\ 15,\ \dots$ là dãy tăng với khoảng cách $4$; \ $90,\ 80,\ 70,\ \dots$ là dãy giảm với khoảng cách $10$.
\item Khi biết số hạng đầu, số hạng cuối và khoảng cách:
\[\text{Số số hạng}=(\text{Số cuối}-\text{Số đầu}):\text{Khoảng cách}+1,\]
\[\text{Tổng}=(\text{Số đầu}+\text{Số cuối})\times\text{Số số hạng}:2.\]
\item Số hạng thứ $n$ của dãy tăng là $\ \text{Số đầu}+(n-1)\times\text{Khoảng cách}$; \ của dãy giảm là $\ \text{Số đầu}-(n-1)\times\text{Khoảng cách}$.
\end{itemize}\end{knowbox}
\gap
\shead{Đếm chữ số -- Lập số từ các chữ số cho trước}
\begin{knowbox}\begin{itemize}[leftmargin=1.4em,itemsep=3pt]
\item Có $9$ số có một chữ số (từ $1$ đến $9$), $90$ số có hai chữ số (từ $10$ đến $99$) và $900$ số có ba chữ số (từ $100$ đến $999$).
\item \textbf{Quy tắc nhân:} để lập một số, ta lần lượt chọn chữ số cho từng hàng. Nếu hàng thứ nhất có $m$ cách chọn, hàng thứ hai có $n$ cách chọn, hàng thứ ba có $p$ cách chọn,\,\dots\ thì có tất cả $m\times n\times p\times\dots$ số.
\item Chữ số $0$ không thể đứng ở hàng cao nhất. Nếu các chữ số phải \emph{khác nhau} thì chữ số đã dùng không được chọn lại.
\item Để tính tổng tất cả các số lập được, ta đếm \emph{mỗi chữ số xuất hiện bao nhiêu lần ở mỗi hàng}.
\end{itemize}\end{knowbox}
\gap
\shead{Cấu tạo số: viết thêm chữ số vào một số}
\begin{knowbox}
Viết số theo cấu tạo thập phân giúp liên hệ số cũ với số mới. Với số có hai chữ số $\nd{ab}$:\par\smallskip
\begin{itemize}[leftmargin=1.4em,itemsep=3pt]
\item viết thêm chữ số $c$ vào bên \textbf{phải}: $\ \nd{abc}=\nd{ab}\times10+c$ (gấp mười lần số cũ rồi cộng thêm $c$);
\item viết thêm chữ số $c$ vào bên \textbf{trái}: $\ \nd{cab}=c\times100+\nd{ab}$;
\item viết thêm chữ số $0$ vào \textbf{giữa}: $\ \nd{a0b}=a\times100+b$.
\end{itemize}
\end{knowbox}
\gap
\begin{exbox}{\ExWord\ 1}
Cho dãy số cách đều $3,\ 7,\ 11,\ 15,\ \dots,\ 199$.
\begin{cols2}
\pa{a} Dãy số có bao nhiêu số hạng? & \pa{b} Tính tổng tất cả các số hạng của dãy.
\end{cols2}
\guide
Khoảng cách giữa hai số hạng liên tiếp là $7-3=4$.\par
\pa{a} Số số hạng: $(199-3):4+1=50$.\par
\pa{b} Tổng các số hạng: $(3+199)\times50:2=5\,050$.
\note{Ghi nhớ: ghép số hạng đầu với số hạng cuối, số hạng thứ hai với số hạng áp cuối,\dots\ mỗi cặp đều có tổng bằng $202$, và có $50:2=25$ cặp.}
\end{exbox}
\gap
\begin{exbox}{\ExWord\ 2}
Cần bao nhiêu chữ số để đánh số trang một cuốn sách dày $234$ trang, bắt đầu từ trang $1$?
\guide
Chia các trang thành ba nhóm theo số chữ số:\par
Từ trang $1$ đến trang $9$: $9$ trang, dùng $9\times1=9$ chữ số.\par
Từ trang $10$ đến trang $99$: $90$ trang, dùng $90\times2=180$ chữ số.\par
Từ trang $100$ đến trang $234$: $234-100+1=135$ trang, dùng $135\times3=405$ chữ số.\par
Tất cả: $9+180+405=594$ chữ số.
\end{exbox}
\gap
\begin{exbox}{\ExWord\ 3}
Từ các chữ số $0,\ 3,\ 5,\ 9$ có thể lập được bao nhiêu số có ba chữ số? Bao nhiêu số có ba chữ số khác nhau?
\guide
Gọi số cần lập là $\nd{abc}$ ($a\ne0$).\par
\textbf{Các chữ số có thể lặp lại:} $a$ có $3$ cách chọn ($3,5,9$), $b$ có $4$ cách, $c$ có $4$ cách.\par Lập được $3\times4\times4=48$ số.\par
\textbf{Các chữ số khác nhau:} $a$ có $3$ cách chọn; $b$ phải khác $a$ nên có $3$ cách (lúc này được chọn $0$); $c$ phải khác cả hai nên có $2$ cách.\par Lập được $3\times3\times2=18$ số.
\end{exbox}
\gap
\begin{exbox}{\ExWord\ 4}
Tính tổng tất cả các số có ba chữ số khác nhau lập được từ $2,\ 3,\ 5,\ 7$.
\guide
Có $4\times3\times2=24$ số.\par Bốn chữ số có vai trò như nhau nên ở mỗi hàng, mỗi chữ số xuất hiện $24:4=6$ lần.\par
Tổng các chữ số ở mỗi hàng là $(2+3+5+7)\times6=102$.\par Vậy tổng cần tìm là:
\[102\times100+102\times10+102\times1=102\times111=11\,322.\]
\end{exbox}
\gap
\begin{exbox}{\ExWord\ 5}
Tìm một số có hai chữ số, biết rằng khi viết thêm chữ số $5$ vào bên phải số đó thì được số mới lớn hơn số đã cho $293$ đơn vị.
\guide
Gọi số cần tìm là $\nd{ab}$.\par Viết thêm $5$ vào bên phải ta được $\nd{ab5}=\nd{ab}\times10+5$.\par
Theo đề bài: $\ \nd{ab}\times10+5=\nd{ab}+293$.\par
Bớt $\nd{ab}$ ở cả hai vế: $\ \nd{ab}\times9+5=293$, \ nên $\ \nd{ab}\times9=288$ và $\nd{ab}=32$.\par
\emph{Thử lại:} $325-32=293$.\par Vậy số cần tìm là $32$.
\end{exbox}
\dividerbar{Luyện tập}
\begin{probbox}{1}{\relax}
Cho dãy số $2,\ 5,\ 8,\ 11,\ 14,\ \dots$
\begin{cols2}
\pa{a} Nêu quy luật và viết tiếp ba số hạng. & \pa{b} Tìm số hạng thứ $100$.\\
\multicolumn{2}{@{}l@{}}{\pa{c} Các số $2\,024$ và $2\,025$ có thuộc dãy không? Nếu có thì đó là số hạng thứ mấy?}
\end{cols2}
\end{probbox}
\probsol{
\pa{a} Mỗi số hạng hơn số hạng đứng liền trước $3$ đơn vị.\par Ba số hạng tiếp theo: $17,\ 20,\ 23$.\par
\pa{b} Số hạng thứ $100$: $2+(100-1)\times3=299$.\par
\pa{c} Mọi số hạng đều có số dư là $2$ khi chia cho $3$.\par Vì $2\,024:3=674$ dư $2$ nên số $2\,024$ thuộc dãy: đó là số hạng thứ $(2\,024-2):3+1=675$.\par Còn $2\,025$ chia hết cho $3$ nên không thuộc dãy.
}
\gap
\begin{probbox}{2}{\relax}
Cho dãy số $4,\ 9,\ 14,\ 19,\ \dots,\ 499$. Tìm số số hạng và tổng của dãy.
\end{probbox}
\probsol{
Khoảng cách: $5$.\par Số số hạng: $(499-4):5+1=100$.\par
Tổng: $(4+499)\times100:2=25\,150$.
\dapso{$100$ số hạng; tổng $25\,150$.}}
\gap
\begin{probbox}{3}{\relax}
Tính tổng của $80$ số tự nhiên liên tiếp bắt đầu từ $25$.
\end{probbox}
\probsol{
Số cuối cùng: $25+(80-1)\times1=104$.\par
Tổng: $(25+104)\times80:2=5\,160$.
\dapso{$5\,160$.}}
\gap
\begin{probbox}{4}{\relax}
Tính tổng $40$ số hạng đầu tiên của dãy số $3,\ 7,\ 11,\ 15,\ \dots$
\end{probbox}
\probsol{
Số hạng thứ $40$: $3+39\times4=159$.\par Tổng: $(3+159)\times40:2=3\,240$.
\dapso{$3\,240$.}}
\gap
\begin{probbox}{5}{\relax}
Tính tổng tất cả các số lẻ có hai chữ số.
\end{probbox}
\probsol{
Các số lẻ có hai chữ số tạo thành dãy $11,\ 13,\ \dots,\ 99$ (khoảng cách $2$).\par
Số số hạng: $(99-11):2+1=45$.\par Tổng: $(11+99)\times45:2=2\,475$.
\dapso{$2\,475$.}}
\gap
\begin{probbox}{6}{\relax}
Cần bao nhiêu chữ số để đánh số trang một cuốn truyện dày $208$ trang, bắt đầu từ trang $1$?
\end{probbox}
\probsol{
Trang $1$--$9$: $9$ chữ số; trang $10$--$99$: $90\times2=180$ chữ số.\par
Trang $100$--$208$: $109$ trang, dùng $109\times3=327$ chữ số.\par
Tất cả: $9+180+327=516$ chữ số.
\dapso{$516$ chữ số.}}
\gap
\begin{probbox}{7}{\relax}
Người ta dùng đúng $435$ chữ số để đánh số trang một cuốn sách, bắt đầu từ trang $1$. Hỏi cuốn sách có bao nhiêu trang?
\end{probbox}
\probsol{
$99$ trang đầu dùng $9+180=189$ chữ số.\par
Số chữ số còn lại: $435-189=246$, dùng cho các trang có ba chữ số: $246:3=82$ trang.\par
Cuốn sách có $99+82=181$ trang.
\dapso{$181$ trang.}}
\gap
\begin{probbox}{8}{\relax}
Viết liên tiếp các số tự nhiên bắt đầu từ $1$ ta được $A=123456789101112\dots$ Chữ số thứ $150$ của $A$ là chữ số nào và thuộc số nào?
\end{probbox}
\probsol{
Các số từ $1$ đến $9$ chiếm $9$ chữ số; $150-9=141$ chữ số còn lại thuộc các số có hai chữ số.\par
$141:2=70$ dư $1$: sau $70$ số có hai chữ số đầu tiên (từ $10$ đến $79$) còn một chữ số nữa, đó là chữ số đầu của số $80$.
\dapso{chữ số $8$, thuộc số $80$.}}
\gap
\begin{probbox}{9}{\relax}
Từ các chữ số $1,\ 4,\ 6,\ 9$ có thể viết được bao nhiêu số tự nhiên
\begin{cols3}
\pa{a} có ba chữ số? & \pa{b} có ba chữ số khác nhau? & \pa{c} có ba chữ số khác nhau và là số chẵn?
\end{cols3}
\end{probbox}
\probsol{
\pa{a} $4\times4\times4=64$. \qquad \pa{b} $4\times3\times2=24$.\par
\pa{c} Chữ số hàng đơn vị là $4$ hoặc $6$: $2$ cách chọn; khi đó hàng trăm có $3$ cách chọn và hàng chục có $2$ cách.\par Lập được $2\times3\times2=12$ số.
}
\gap
\begin{probbox}{10}{\relax}
Từ các chữ số $0,\ 2,\ 5,\ 7$ có thể viết được bao nhiêu số có ba chữ số? Bao nhiêu số có ba chữ số khác nhau?
\end{probbox}
\probsol{
Chữ số hàng trăm không thể là $0$ nên có $3$ cách chọn.\par
Các chữ số có thể lặp lại: $3\times4\times4=48$.\par Các chữ số khác nhau: $3\times3\times2=18$.
\dapso{$48$ và $18$.}}
\gap
\begin{probbox}{11}{\relax}
Từ các chữ số $0,\ 1,\ 3,\ 6,\ 8$ có thể viết được bao nhiêu số có bốn chữ số khác nhau? Trong đó có bao nhiêu số chẵn?
\end{probbox}
\probsol{
\textbf{Tất cả các số:} chữ số hàng nghìn có $4$ cách chọn (khác $0$), tiếp theo là $4$, $3$, $2$ cách chọn: $4\times4\times3\times2=96$.\par
\textbf{Các số chẵn} (tận cùng là $0$, $6$ hoặc $8$):\par
-- Tận cùng là $0$: $4\times3\times2=24$.\par
-- Tận cùng là $6$ (hoặc $8$): chữ số hàng nghìn có $3$ cách chọn (khác $0$, khác chữ số tận cùng), tiếp theo là $3$ và $2$: $3\times3\times2=18$.\par Cả hai trường hợp được $36$.\par
Số các số chẵn: $24+36=60$.
\dapso{$96$ số, trong đó có $60$ số chẵn.}}
\gap
\begin{probbox}{12}{\relax}
Tính tổng tất cả các số có ba chữ số khác nhau lập được từ $1,\ 3,\ 5,\ 7$.
\end{probbox}
\probsol{
Có $4\times3\times2=24$ số; ở mỗi hàng, mỗi chữ số xuất hiện $24:4=6$ lần.\par
Tổng: $(1+3+5+7)\times6\times111=16\times6\times111=10\,656$.
\dapso{$10\,656$.}}
\gap
\begin{probbox}{13}{\relax}
Có bao nhiêu số có ba chữ số chứa đúng một chữ số $5$?
\end{probbox}
\probsol{
-- $5$ ở hàng trăm: hàng chục và hàng đơn vị mỗi hàng có $9$ cách chọn (khác $5$): $9\times9=81$.\par
-- $5$ ở hàng chục: hàng trăm có $8$ cách chọn (khác $0$, khác $5$), hàng đơn vị có $9$ cách: $8\times9=72$.\par
-- $5$ ở hàng đơn vị: tương tự, $8\times9=72$.\par
Tất cả: $81+72+72=225$.
\dapso{$225$ số.}}
\gap
\begin{probbox}{14}{\relax}
Có bao nhiêu số có bốn chữ số khác nhau mà chữ số hàng nghìn là $3$?
\end{probbox}
\probsol{
Chữ số hàng nghìn cố định là $3$.\par Hàng trăm có $9$ cách chọn (bất kì chữ số nào trừ $3$), hàng chục có $8$ cách, hàng đơn vị có $7$ cách: $9\times8\times7=504$.
\dapso{$504$ số.}}
\gap
\begin{probbox}{15}{\relax}
Tìm một số có hai chữ số, biết rằng khi viết thêm chữ số $4$ vào bên phải số đó thì được số mới lớn hơn số đã cho $571$ đơn vị.
\end{probbox}
\probsol{
Gọi số cần tìm là $\nd{ab}$.\par Ta có $\nd{ab}\times10+4=\nd{ab}+571$, nên $\nd{ab}\times9=567$ và $\nd{ab}=63$.\par
\emph{Thử lại:} $634-63=571$.
\dapso{$63$.}}
\gap
\begin{probbox}{16}{\relax}
Tìm một số có ba chữ số, biết rằng khi viết thêm chữ số $9$ vào bên phải số đó thì được số mới lớn hơn số đã cho $1\,350$ đơn vị.
\end{probbox}
\probsol{
Gọi số cần tìm là $x$.\par Số mới là $x\times10+9$, nên $x\times10+9=x+1\,350$, suy ra $x\times9=1\,341$ và $x=149$.\par
\emph{Thử lại:} $1\,499-149=1\,350$.
\dapso{$149$.}}
\gap
\begin{probbox}{17}{\relax}
Tìm một số có ba chữ số, biết rằng khi viết thêm chữ số $4$ vào bên trái số đó thì được số mới gấp $17$ lần số đã cho.
\end{probbox}
\probsol{
Gọi số cần tìm là $x$.\par Số mới là $4\,000+x$, nên $4\,000+x=17\times x$, suy ra $16\times x=4\,000$ và $x=250$.\par
\emph{Thử lại:} $4\,250=17\times250$.
\dapso{$250$.}}
\gap
\begin{probbox}{18}{\relax}
Tìm một số có hai chữ số, biết rằng khi viết thêm chữ số $0$ vào giữa hai chữ số của nó thì được số mới gấp $7$ lần số đã cho.
\end{probbox}
\probsol{
Gọi số cần tìm là $\nd{ab}$.\par Ta có $\nd{a0b}=7\times\nd{ab}$, tức là\par $a\times100+b=a\times70+b\times7$.\par
Suy ra $a\times30=b\times6$, tức là\par $b=a\times5$.\par Vì $b$ là chữ số nên $a=1$ và $b=5$.\par
\emph{Thử lại:} $105=7\times15$.
\dapso{$15$.}}
\gap
\begin{probbox}{19}{\relax}
Cho dãy số $5,\ 10,\ 15,\ \dots,\ 500$. Dãy có bao nhiêu số hạng? Tính tổng các số hạng đó.
\end{probbox}
\probsol{
Số số hạng: $(500-5):5+1=100$.\par Tổng: $(5+500)\times100:2=25\,250$.
\dapso{$100$ số hạng; $25\,250$.}}
\gap
\begin{probbox}{20}{\relax}
Cho dãy số $1,\ 4,\ 7,\ 10,\ \dots$ \ \pa{a} Tìm số hạng thứ $50$. \ \pa{b} Số $2\,026$ có thuộc dãy không? Nếu có thì đó là số hạng thứ mấy?
\end{probbox}
\probsol{
\pa{a} Số hạng thứ $50$: $1+(50-1)\times3=148$.\par
\pa{b} Mọi số hạng đều có số dư là $1$ khi chia cho $3$; $2\,026=3\times675+1$, nên số này thuộc dãy.\par Đó là số hạng thứ $(2\,026-1):3+1=676$.
\dapso{\pa{a} $148$;\par \pa{b} có, là số hạng thứ $676$.}}
\gap
\begin{probbox}{21}{\relax}
Tính tổng các số chẵn từ $2$ đến $200$.
\end{probbox}
\probsol{
Số số hạng: $(200-2):2+1=100$.\par Tổng: $(2+200)\times100:2=10\,100$.
\dapso{$10\,100$.}}
\gap
\begin{probbox}{22}{\relax}
Tính tổng tất cả các số có ba chữ số chia hết cho $10$.
\end{probbox}
\probsol{
Dãy $100,\ 110,\ \dots,\ 990$ có $(990-100):10+1=90$ số hạng.\par Tổng: $(100+990)\times90:2=49\,050$.
\dapso{$49\,050$.}}
\gap
\begin{probbox}{23}{\relax}
Có bao nhiêu số có hai chữ số chia hết cho $3$?
\end{probbox}
\probsol{
Đó là các số $12,\ 15,\ \dots,\ 99$.\par Số lượng: $(99-12):3+1=30$.
\dapso{$30$ số.}}
\gap
\begin{probbox}{24}{\relax}
Cần bao nhiêu chữ số để đánh số trang một cuốn sách dày $120$ trang, bắt đầu từ trang $1$?
\end{probbox}
\probsol{
Trang $1$ -- $9$: $9$ chữ số.\par Trang $10$ -- $99$: $90\times2=180$ chữ số.\par Trang $100$ -- $120$: $21\times3=63$ chữ số.\par Tất cả: $9+180+63=252$ chữ số.
\dapso{$252$ chữ số.}}
\gap
\begin{probbox}{25}{\relax}
Một con phố có $75$ ngôi nhà được đánh số từ $1$ đến $75$. Hỏi phải dùng bao nhiêu chữ số để đánh số các nhà đó?
\end{probbox}
\probsol{
Số có một chữ số: $9$ số, dùng $9$ chữ số.\par Số có hai chữ số: $75-9=66$ số, dùng $66\times2=132$ chữ số.\par Tất cả: $9+132=141$ chữ số.
\dapso{$141$ chữ số.}}
\gap
\begin{probbox}{26}{\relax}
Từ các chữ số $2,\ 4,\ 6$, hãy viết tất cả các số có ba chữ số khác nhau rồi tính tổng các số đó.
\end{probbox}
\probsol{
Các số đó là: $246,\ 264,\ 426,\ 462,\ 624,\ 642$ ($6$ số).\par Mỗi chữ số xuất hiện ở mỗi hàng đúng $2$ lần.\par Tổng: $(2+4+6)\times2\times111=2\,664$.
\dapso{$2\,664$.}}
\gap
\begin{probbox}{27}{\relax}
Từ $1,\ 2,\ 3,\ 4$ có thể viết được bao nhiêu số có hai chữ số khác nhau?
\end{probbox}
\probsol{
Chữ số hàng chục: $4$ cách chọn.\par Chữ số hàng đơn vị: $3$ cách chọn (khác chữ số hàng chục).\par Số lượng: $4\times3=12$.
\dapso{$12$ số.}}
\gap
\begin{probbox}{28}{\relax}
Tìm một số có hai chữ số, biết rằng số đó tăng thêm $362$ đơn vị khi viết thêm chữ số $2$ vào bên phải.
\end{probbox}
\probsol{
Số mới gấp $10$ lần số cũ rồi cộng thêm $2$.\par Vậy $9$ lần số cũ bằng $362-2=360$.\par Số cần tìm là $360:9=40$.
\dapso{$40$.}}
\gap
\begin{probbox}{29}{\relax}
Tìm một số có ba chữ số, biết rằng khi viết thêm chữ số $1$ vào bên trái số đó thì được số mới gấp $9$ lần số đã cho.
\end{probbox}
\probsol{
Viết thêm $1$ vào bên trái thì số đó tăng thêm $1\,000$.\par Số mới gấp $9$ lần số đã cho, nên $1\,000$ gấp $9-1=8$ lần số đó.\par Số cần tìm là $1\,000:8=125$.
\dapso{$125$.}}
\gap
\begin{probbox}{30}{\relax}
Một số tự nhiên có chữ số tận cùng là $7$. Nếu xoá chữ số $7$ đó thì số đã cho giảm đi $1\,789$ đơn vị. Tìm số đó.
\end{probbox}
\probsol{
Gọi số sau khi xoá là A; số ban đầu là $10\times A+7$.\par Hiệu là $9\times A+7=1\,789$, nên $9\times A=1\,782$ và $A=198$.\par Số cần tìm là $1\,987$.
\dapso{$1\,987$.}}
\dividerbar{Thử thách -- Nâng cao}
\begin{probbox}{31}{\stars{1}}
Từ $1,\ 2,\ 5,\ 8,\ 9$ có thể viết được bao nhiêu số có bốn chữ số khác nhau mà hai chữ số $2$ và $5$ đứng cạnh nhau?
\end{probbox}
\probsol{
Coi $2$ và $5$ là một ``khối'', viết là $(25)$ hoặc $(52)$: $2$ cách.\par
Chọn thêm $2$ chữ số từ $1,\ 8,\ 9$: $3$ cách (bỏ đi một chữ số).\par
Sắp xếp khối và hai chữ số đó (ba ``phần tử''): $3\times2\times1=6$ cách.\par
Tất cả: $2\times3\times6=36$.
\dapso{$36$ số.}}
\gap
\begin{probbox}{32}{\stars{1}}
Một số có ba chữ số có chữ số hàng trăm là $7$. Nếu chuyển chữ số $7$ này sang tận cùng bên phải thì được số mới bé hơn số ban đầu $324$ đơn vị. Tìm số đó.
\end{probbox}
\probsol{
Gọi số cần tìm là $\nd{7bc}=700+\nd{bc}$.\par Số mới là $\nd{bc7}=\nd{bc}\times10+7$.\par
Ta có $700+\nd{bc}-(\nd{bc}\times10+7)=324$, tức là\par $693-\nd{bc}\times9=324$.\par
Suy ra $\nd{bc}\times9=369$ và $\nd{bc}=41$.\par Số cần tìm là $741$ (\emph{thử lại:} $741-417=324$).
\dapso{$741$.}}
\gap
\begin{probbox}{33}{\stars{1}}
Tìm một số có hai chữ số, biết rằng khi viết thêm số $12$ vào giữa hai chữ số của nó thì được số mới gấp $76$ lần số đã cho.
\end{probbox}
\probsol{
Gọi số cần tìm là $\nd{ab}$.\par Số mới là $\nd{a12b}=a\times1\,000+120+b$.\par
Ta có $a\times1\,000+120+b=76\times(a\times10+b)=a\times760+b\times76$.\par
Suy ra $a\times240+120=b\times75$; chia cả hai vế cho $15$: $a\times16+8=b\times5$.\par
Vế phải không quá $45$ nên $a\le2$; với $a=1$ được $24$, không chia hết cho $5$; với $a=2$ được $40=b\times5$, $b=8$.\par
\emph{Thử lại:} $2\,128=76\times28$.
\dapso{$28$.}}
\gap
\begin{probbox}{34}{\stars{1}}
Viết liên tiếp các số chẵn bắt đầu từ $2$ ta được $\ 2468101214\dots$ Chữ số thứ $200$ là chữ số nào?
\end{probbox}
\probsol{
Các số chẵn có một chữ số ($2,4,6,8$) dùng $4$ chữ số.\par Các số chẵn có hai chữ số (từ $10$ đến $98$) gồm $45$ số, dùng $90$ chữ số.\par Đến đây đã có: $94$ chữ số.\par
Còn lại: $200-94=106$ chữ số; $106:3=35$ dư $1$.\par Sau $35$ số chẵn có ba chữ số (từ $100$ đến $168$) là chữ số đầu của số $170$.
\dapso{chữ số $1$ (của số $170$).}}
\gap
\begin{probbox}{35}{\stars{1}}
Một dãy số cách đều tăng dần có $30$ số hạng, số hạng đầu là $9$ và tổng bằng $2\,445$. Tìm khoảng cách và số hạng cuối của dãy.
\end{probbox}
\probsol{
Từ công thức tính tổng: $(9+\text{số cuối})\times30:2=2\,445$, nên $9+\text{số cuối}=2\,445\times2:30=163$ và số hạng cuối là $154$.\par
Từ số hạng đầu đến số hạng cuối có $29$ khoảng cách: $(154-9):29=5$.
\dapso{khoảng cách $5$, số hạng cuối $154$.}}
\gap
\begin{probbox}{36}{\stars{2}}
Tìm số tự nhiên lớn nhất có các chữ số khác nhau mà tổng các chữ số bằng $20$.
\end{probbox}
\probsol{
Số càng có nhiều chữ số thì càng lớn.\par Bảy chữ số khác nhau có tổng ít nhất là $0+1+\dots+6=21>20$, nên số cần tìm có nhiều nhất $6$ chữ số.\par
Để số lớn nhất, chữ số hàng cao nhất phải lớn nhất có thể: chọn $9$.\par Năm chữ số khác nhau còn lại có tổng bằng $11$; năm chữ số nhỏ nhất có tổng $0+1+2+3+4=10$, nên chúng chỉ có thể là $0,1,2,3,5$.\par
Xếp các chữ số theo thứ tự giảm dần: $953\,210$.
\dapso{$953\,210$.}}
\gap
\begin{probbox}{37}{\stars{1}}
An viết các số tự nhiên từ $1$ đến $100$ lên bảng. Hỏi An đã viết bao nhiêu lần chữ số $1$?
\end{probbox}
\probsol{
-- Ở hàng đơn vị: $1,\ 11,\ 21,\ \dots,\ 91$: $10$ lần.\par
-- Ở hàng chục: $10,\ 11,\ \dots,\ 19$: $10$ lần.\par
-- Ở hàng trăm: số $100$: một lần.\par
Tất cả: $10+10+1=21$ lần.
\dapso{$21$ lần.}}
\gap
\begin{probbox}{38}{\stars{1}}
Tìm số tự nhiên $n$ sao cho $1+2+3+\dots+n=210$.
\end{probbox}
\probsol{
Tổng bằng $(1+n)\times n:2$, nên $n\times(n+1)=420$.\par Hai số tự nhiên liên tiếp có tích bằng $420=20\times21$.\par Vậy $n=20$.
\dapso{$n=20$.}}
\gap
\begin{probbox}{39}{\stars{1}}
Trong dãy $1,\ 3,\ 6,\ 10,\ 15,\ \dots$ mỗi số hạng bằng số hạng đứng trước cộng với số thứ tự của nó. Tìm số hạng thứ $20$.
\end{probbox}
\probsol{
Số hạng thứ $n$ là $1+2+\dots+n$.\par Số hạng thứ $20$: $1+2+\dots+20=(1+20)\times20:2=210$.
\dapso{$210$.}}
\gap
\begin{probbox}{40}{\stars{1}}
Tính $S=1+2+4+8+16+\dots+512$ (mỗi số hạng gấp đôi số hạng đứng trước).
\end{probbox}
\probsol{
$S\times2=2+4+8+\dots+512+1\,024$.\par Trừ $S$ khỏi $S\times2$, các số hạng giống nhau triệt tiêu: $S=1\,024-1=1\,023$.
\dapso{$1\,023$.}}
\gap
\begin{probbox}{41}{\stars{1}}
Viết liên tiếp các số từ $1$ đến $50$ ta được một số rất dài. \ \pa{a} Số đó có bao nhiêu chữ số? \ \pa{b} Chữ số thứ $50$ kể từ trái sang là chữ số nào?
\end{probbox}
\probsol{
\pa{a} $9+41\times2=91$ chữ số.\par
\pa{b} $9$ chữ số đầu là $1$ -- $9$; $50-9=41$ chữ số còn lại thuộc các số có hai chữ số.\par $41=20\times2+1$, nên viết đủ $20$ số ($10$ -- $29$) và đến chữ số đầu của số $30$.\par Chữ số thứ $50$ là $3$.
\dapso{\pa{a} $91$;\par \pa{b} $3$.}}
\gap
\begin{probbox}{42}{\stars{1}}
Có bao nhiêu số có ba chữ số mà tổng các chữ số bằng $5$?
\end{probbox}
\probsol{
Xét theo chữ số hàng trăm:\par
Chữ số hàng trăm là $1$: hai chữ số còn lại có tổng bằng $4$, được $5$ số ($104,\ 113,\ 122,\ 131,\ 140$).\par Chữ số hàng trăm là $2,\ 3,\ 4,\ 5$: lần lượt có $4,\ 3,\ 2,\ 1$ số.\par Tất cả: $5+4+3+2+1=15$.
\dapso{$15$ số.}}
\gap
\begin{probbox}{43}{\stars{1}}
Có bao nhiêu số có ba chữ số chia hết cho $5$?
\end{probbox}
\probsol{
Đó là các số $100,\ 105,\ \dots,\ 995$.\par Số lượng: $(995-100):5+1=180$.
\dapso{$180$ số.}}
\gap
\begin{probbox}{44}{\stars{1}}
Từ $0,\ 1,\ 2,\ 3,\ 4$ có thể viết được bao nhiêu số chẵn có ba chữ số khác nhau?
\end{probbox}
\probsol{
\emph{Chữ số hàng đơn vị là $0$:} hàng trăm có $4$ cách chọn, hàng chục có $3$ cách: $12$ số.\par
\emph{Chữ số hàng đơn vị là $2$ hoặc $4$:} hàng trăm có $3$ cách chọn (khác $0$, khác chữ số hàng đơn vị), hàng chục có $3$ cách: $2\times3\times3=18$.\par Tất cả: $12+18=30$.
\dapso{$30$ số.}}
\gap
\begin{probbox}{45}{\stars{1}}
Tìm tất cả các số có hai chữ số gấp $4$ lần tổng các chữ số của nó.
\end{probbox}
\probsol{
$\nd{ab}=a\times10+b$ và $4\times(a+b)=a\times4+b\times4$.\par Hai biểu thức bằng nhau khi $a\times6=b\times3$, tức là\par $b=a\times2$.\par Các số đó là: $12,\ 24,\ 36,\ 48$.
\dapso{$12;\ 24;\ 36;\ 48$.}}
\gap
\begin{probbox}{46}{\stars{1}}
Một số có hai chữ số có tổng các chữ số bằng $10$. Nếu đổi chỗ hai chữ số của nó thì được số mới lớn hơn số đã cho $36$ đơn vị. Tìm số đó.
\end{probbox}
\probsol{
Khi đổi chỗ hai chữ số, số đó thay đổi $9$ lần hiệu hai chữ số.\par Hiệu hai chữ số là $36:9=4$, chữ số lớn hơn ở hàng đơn vị.\par Chữ số hàng chục: $(10-4):2=3$; chữ số hàng đơn vị: $7$.
\dapso{$37$.}}
\gap
\begin{probbox}{47}{\stars{1}}
Người ta dùng $1\,092$ chữ số để đánh số trang một cuốn sách, bắt đầu từ trang $1$. Hỏi cuốn sách có bao nhiêu trang?
\end{probbox}
\probsol{
Trang $1$ -- $99$ dùng $9+180=189$ chữ số.\par $1\,092-189=903$ chữ số còn lại dùng để đánh $903:3=301$ trang có ba chữ số.\par Số trang: $99+301=400$.
\dapso{$400$ trang.}}
\gap
\begin{probbox}{48}{\stars{2}}
Khi viết các số từ $1$ đến $200$, chữ số $2$ được viết bao nhiêu lần?
\end{probbox}
\probsol{
Ở hàng đơn vị: $2,\ 12,\ \dots,\ 192$, $20$ lần.\par Ở hàng chục: $20$ -- $29$ và $120$ -- $129$, $20$ lần.\par Ở hàng trăm: chỉ có số $200$, một lần.\par Tất cả: $20+20+1=41$ lần.
\dapso{$41$ lần.}}
\gap
\begin{probbox}{49}{\stars{2}}
Tìm số tự nhiên nhỏ nhất có tổng các chữ số bằng $30$.
\end{probbox}
\probsol{
Để số nhỏ nhất thì số đó phải có ít chữ số nhất; mỗi chữ số không quá $9$, mà $9\times3=27<30$, nên số đó có ít nhất $4$ chữ số.\par Chữ số hàng cao nhất phải nhỏ nhất có thể, theo sau là ba chữ số $9$: $30-27=3$.\par Số cần tìm là $3\,999$.
\dapso{$3\,999$.}}
\gap
\begin{probbox}{50}{\stars{2}}
Năm số tự nhiên liên tiếp có tổng bằng $2\,025$. Tìm năm số đó.
\end{probbox}
\probsol{
Số ở giữa bằng trung bình cộng: $2\,025:5=405$.\par Các số đó là: $403,\ 404,\ 405,\ 406,\ 407$.
\dapso{$403;\ 404;\ 405;\ 406;\ 407$.}}
\clearpage
\titlehead
\begingroup\setbox0=\vbox{\section*{Phân số và hỗn số}}\endgroup
\chapterhead{2}{Phân số và hỗn số}

\shead{Phân số -- Tính chất cơ bản}
\begin{knowbox}\begin{itemize}[leftmargin=1.4em,itemsep=3pt]
\item Phân số có dạng $\dfrac{a}{b}$, trong đó $a,\ b$ là các số tự nhiên và $b\ne0$; $a$ là \textbf{tử số}, $b$ là \textbf{mẫu số}. Thương $a:b$ là giá trị của phân số $\dfrac{a}{b}$.
\item \textbf{Tính chất cơ bản:} nếu nhân (hoặc chia) cả tử số và mẫu số với cùng một số khác không thì được một phân số bằng phân số đã cho:
\[\frac{a}{b}=\frac{a\times n}{b\times n}=\frac{a:m}{b:m}\qquad(n,\ m\ne0).\]
\item Nhờ đó ta có thể \textbf{rút gọn} phân số (đến phân số tối giản) và \textbf{quy đồng mẫu số} các phân số.
\item Hai phân số bằng nhau, $\ \dfrac{a}{b}=\dfrac{c}{d}$, khi và chỉ khi $\ a\times d=b\times c$ (nhân chéo).
\end{itemize}\end{knowbox}
\gap
\shead{Hỗn số}
\begin{knowbox}
Khi tử số của $\dfrac{a}{b}$ lớn hơn mẫu số, ta chia $a$ cho $b$ được thương $q$ và số dư $r$ ($0\le r<b$). Khi đó
\[\frac{a}{b}=q\frac{r}{b}\quad\text{(đọc là ``$q$ và $\tfrac{r}{b}$'')},\qquad\text{ví dụ }\ \frac{17}{5}=3\frac{2}{5}.\]
$q\dfrac{r}{b}$ là một \textbf{hỗn số}: gồm \emph{phần nguyên} $q$ và \emph{phần phân số} $\dfrac{r}{b}$ (bé hơn $1$). Ngược lại, $\ q\dfrac{r}{b}=\dfrac{q\times b+r}{b}$.
\end{knowbox}
\gap
\shead{Tính chất của các phép tính -- Tính nhanh}
\begin{knowbox}
Với các phân số (hoặc hỗn số) $A,\ B,\ C$:\par\smallskip
\begin{cols2}
\textbf{Giao hoán:} $A+B=B+A$;\ \ $A\times B=B\times A$ & \textbf{Kết hợp:} $(A+B)+C=A+(B+C)$\\
\textbf{Trừ một số cho một tổng:} $A-(B+C)=A-B-C$ & \textbf{Phân phối:} $A\times(B\pm C)=A\times B\pm A\times C$\\
\textbf{Trừ một số cho một hiệu:} $A-(B-C)=A-B+C$ & \textbf{Chia một số cho một tích:} $A:(B\times C)=A:B:C$
\end{cols2}
\smallskip Để tính nhanh, ta \emph{nhóm các số hạng có tổng (hoặc hiệu) là số tự nhiên}, hoặc \emph{đưa thừa số chung ra ngoài}.
\end{knowbox}
\gap
\shead{Tổng các phân số viết theo quy luật}
\begin{knowbox}\begin{itemize}[leftmargin=1.4em,itemsep=3pt]
\item \textbf{Tách một phân số thành hiệu:} với $b>a$,
\[\frac{b-a}{a\times b}=\frac{1}{a}-\frac{1}{b},\qquad\text{ví dụ }\ \frac{1}{2\times3}=\frac12-\frac13;\quad \frac{3}{4\times7}=\frac14-\frac17.\]
Khi cộng các phân số như vậy, các số hạng ở giữa \emph{triệt tiêu nhau}, chỉ còn lại số hạng đầu và số hạng cuối.
\item \textbf{Tổng có số hạng sau gấp số hạng trước một số lần không đổi} (mỗi số hạng bằng số hạng đứng trước nhân với cùng một số), ví dụ $\dfrac12+\dfrac14+\dfrac18+\dots$: nhân tổng với số đó rồi trừ đi tổng ban đầu.
\end{itemize}\end{knowbox}
\gap
\begin{exbox}{\ExWord\ 1}
Tính bằng cách thuận tiện nhất:
\begin{cols3}
$A=\dfrac35+2\dfrac14+\dfrac25+\dfrac34$ & $B=4\dfrac27-\left(1\dfrac38+\dfrac27\right)$ & $C=\dfrac74\times3\dfrac13-\dfrac74\times1\dfrac13$
\end{cols3}
\guide
\pa{A} Nhóm các số hạng có tổng là số tự nhiên: $A=\left(\dfrac35+\dfrac25\right)+\left(2\dfrac14+\dfrac34\right)=1+3=4.$\par
\pa{B} Trừ một số cho một tổng: $B=4\dfrac27-\dfrac27-1\dfrac38=4-1\dfrac38=2\dfrac58.$\par
\pa{C} Đưa thừa số chung ra ngoài: $C=\dfrac74\times\left(3\dfrac13-1\dfrac13\right)=\dfrac74\times2=\dfrac72.$
\end{exbox}
\gap
\begin{exbox}{\ExWord\ 2}
Tính:
\begin{cols2}
$A=\dfrac{1}{1\times2}+\dfrac{1}{2\times3}+\dfrac{1}{3\times4}+\dots+\dfrac{1}{9\times10}$ & $B=\dfrac16+\dfrac1{12}+\dfrac1{20}+\dots+\dfrac1{110}$
\end{cols2}
\guide
\pa{A} Tách mỗi phân số: $\dfrac{1}{1\times2}=\dfrac11-\dfrac12$, \ $\dfrac{1}{2\times3}=\dfrac12-\dfrac13$, \dots
\[A=1-\frac12+\frac12-\frac13+\dots+\frac19-\frac1{10}=1-\frac1{10}=\frac9{10}.\]
\pa{B} Nhận xét: $6=2\times3$, $12=3\times4$, $20=4\times5$,\,\dots, $110=10\times11$.\par Do đó
\[B=\frac12-\frac13+\frac13-\frac14+\dots+\frac1{10}-\frac1{11}=\frac12-\frac1{11}=\frac9{22}.\]
\end{exbox}
\gap
\begin{exbox}{\ExWord\ 3}
Tính $D=\dfrac13+\dfrac19+\dfrac1{27}+\dfrac1{81}+\dfrac1{243}$.
\guide
Mỗi số hạng bằng số hạng đứng trước nhân với $\dfrac13$.\par Nhân $D$ với $3$:
\[3\times D=1+\frac13+\frac19+\frac1{27}+\frac1{81}.\]
Trừ vế với vế, các số hạng giống nhau triệt tiêu: $\ 3\times D-D=1-\dfrac1{243}$, nên $\ 2\times D=\dfrac{242}{243}$.\par
Vậy $D=\dfrac{242}{243}:2=\dfrac{121}{243}$.
\end{exbox}
\gap
\begin{exbox}{\ExWord\ 4}
Tính tích $\ P=\left(1-\dfrac12\right)\times\left(1-\dfrac13\right)\times\left(1-\dfrac14\right)\times\dots\times\left(1-\dfrac1{20}\right)$.
\guide
Tính từng thừa số rồi rút gọn chéo (mỗi tử số rút gọn với mẫu số của phân số đứng trước):
\[P=\frac12\times\frac23\times\frac34\times\dots\times\frac{19}{20}=\frac{1\times2\times3\times\dots\times19}{2\times3\times4\times\dots\times20}=\frac1{20}.\]
\end{exbox}
\dividerbar{Luyện tập}
\begin{probbox}{1}{\relax}
Rút gọn thành phân số tối giản:
\[\frac{18}{24};\qquad\frac{45}{60};\qquad\frac{56}{98};\qquad\frac{125}{1\,000}.\]
\end{probbox}
\probsol{
Chia cả tử số và mẫu số cho ước chung lớn nhất của chúng:
\[\frac{18}{24}=\frac{18:6}{24:6}=\frac34;\quad\frac{45}{60}=\frac{45:15}{60:15}=\frac34;\quad\frac{56}{98}=\frac{56:14}{98:14}=\frac47;\quad\frac{125}{1\,000}=\frac{125:125}{1\,000:125}=\frac18.\]
}
\gap
\begin{probbox}{2}{\relax}
Điền dấu $<,\ >,\ =$ và giải thích:
\begin{cols3}
\pa{a} $\dfrac58\ \square\ \dfrac7{12}$ & \pa{b} $\dfrac{17}{19}\ \square\ \dfrac{17}{23}$ & \pa{c} $\dfrac{2\,024}{2\,025}\ \square\ \dfrac{2\,025}{2\,026}$
\end{cols3}
\end{probbox}
\probsol{
\pa{a} Quy đồng mẫu số: $\dfrac58=\dfrac{15}{24}$, $\dfrac7{12}=\dfrac{14}{24}$.\par Vì $15>14$ nên $\dfrac58>\dfrac7{12}$.\par
\pa{b} Hai phân số có cùng tử số: phân số nào có mẫu số bé hơn thì lớn hơn, nên $\dfrac{17}{19}>\dfrac{17}{23}$.\par
\pa{c} So sánh ``phần bù'' đến $1$: $1-\dfrac{2\,024}{2\,025}=\dfrac1{2\,025}$ lớn hơn $1-\dfrac{2\,025}{2\,026}=\dfrac1{2\,026}$.\par Phân số nào có phần bù lớn hơn thì bé hơn, nên $\dfrac{2\,024}{2\,025}<\dfrac{2\,025}{2\,026}$.
}
\gap
\begin{probbox}{3}{\relax}
Tính:
\begin{cols3}
$A=\dfrac79-\dfrac4{15}+\dfrac16$ & $B=\dfrac{38}{15}-\left(\dfrac23+\dfrac8{15}\right)$ & $C=\dfrac{16}{11}-\dfrac56+\dfrac6{11}-\dfrac16$
\end{cols3}
\end{probbox}
\probsol{
\pa{A} Mẫu số chung $90$: $A=\dfrac{70}{90}-\dfrac{24}{90}+\dfrac{15}{90}=\dfrac{61}{90}$.\par
\pa{B} $B=\dfrac{38}{15}-\dfrac8{15}-\dfrac23=2-\dfrac23=\dfrac43$.\par
\pa{C} $C=\left(\dfrac{16}{11}+\dfrac6{11}\right)-\left(\dfrac56+\dfrac16\right)=2-1=1$.
}
\gap
\begin{probbox}{4}{\relax}
Tính:
\begin{cols3}
$A=2\dfrac35-1\dfrac13+1\dfrac7{15}$ & $B=5\dfrac14+2\dfrac56-3\dfrac1{12}$ & $C=6\dfrac27-\left(2\dfrac13+1\dfrac27\right)$
\end{cols3}
\end{probbox}
\probsol{
\pa{A} $A=\dfrac{39}{15}-\dfrac{20}{15}+\dfrac{22}{15}=\dfrac{41}{15}=2\dfrac{11}{15}$.\par
\pa{B} $B=\dfrac{63}{12}+\dfrac{34}{12}-\dfrac{37}{12}=\dfrac{60}{12}=5$.\par
\pa{C} $C=\left(6\dfrac27-1\dfrac27\right)-2\dfrac13=5-2\dfrac13=2\dfrac23$.
}
\gap
\begin{probbox}{5}{\relax}
Tính bằng cách thuận tiện nhất:
\begin{cols2}
$A=\dfrac57\times\dfrac38+\dfrac57\times\dfrac58$ & $B=\dfrac34\times\dfrac{12}5\times\dfrac59\times\dfrac43$\\
$C=\dfrac79\times\dfrac56-\dfrac79\times\dfrac13$ & $D=2\dfrac13:1\dfrac16$
\end{cols2}
\end{probbox}
\probsol{
\pa{A} $A=\dfrac57\times\left(\dfrac38+\dfrac58\right)=\dfrac57\times1=\dfrac57$.\par
\pa{B} $B=\left(\dfrac34\times\dfrac43\right)\times\left(\dfrac{12}5\times\dfrac59\right)=1\times\dfrac43=\dfrac43$.\par
\pa{C} $C=\dfrac79\times\left(\dfrac56-\dfrac13\right)=\dfrac79\times\dfrac12=\dfrac7{18}$.\par
\pa{D} $D=\dfrac73:\dfrac76=\dfrac73\times\dfrac67=2$.
}
\gap
\begin{probbox}{6}{\relax}
Tìm $y$:
\begin{cols3}
\pa{a} $y\times\dfrac35=\dfrac9{10}$ & \pa{b} $y+\dfrac56=2\dfrac14$ & \pa{c} $\dfrac47:y=\dfrac23$\\
\pa{d} $\left(y-\dfrac12\right)\times\dfrac45=\dfrac25$ & \pa{e} $y\times1\dfrac14+y\times\dfrac34=6$ & \pa{g} $y:\dfrac23+y:\dfrac13=18$
\end{cols3}
\end{probbox}
\probsol{
\pa{a} $y=\dfrac9{10}:\dfrac35=\dfrac9{10}\times\dfrac53=\dfrac32$. \qquad \pa{b} $y=\dfrac94-\dfrac56=\dfrac{27}{12}-\dfrac{10}{12}=\dfrac{17}{12}$.\par
\pa{c} $y=\dfrac47:\dfrac23=\dfrac47\times\dfrac32=\dfrac67$. \qquad \pa{d} $y-\dfrac12=\dfrac25:\dfrac45=\dfrac12$, \ $y=1$.\par
\pa{e} $y\times\left(1\dfrac14+\dfrac34\right)=6$, tức là\par $y\times2=6$, \ $y=3$.\par
\pa{g} $y\times\dfrac32+y\times3=18$, tức là\par $y\times\dfrac92=18$, \ $y=4$.
}
\gap
\begin{probbox}{7}{\relax}
Tìm số tự nhiên $y$:
\begin{cols2}
\pa{a} $\dfrac{y-1}{4}=\dfrac{15}{20}$ & \pa{b} $\dfrac8{12}=\dfrac{10}{y+3}$\\
\pa{c} $\dfrac14<\dfrac2y<\dfrac12$ & \pa{d} $\dfrac37<\dfrac{y}{14}<\dfrac57$
\end{cols2}
\end{probbox}
\probsol{
\pa{a} $\dfrac{15}{20}=\dfrac34$, nên $y-1=3$, $y=4$. \qquad \pa{b} $\dfrac8{12}=\dfrac23=\dfrac{10}{15}$, nên $y+3=15$, $y=12$.\par
\pa{c} Đưa về cùng tử số $2$: $\dfrac28<\dfrac2y<\dfrac24$, nên $4<y<8$: $\ y\in\{5;6;7\}$.\par
\pa{d} Đưa về cùng mẫu số $14$: $\dfrac6{14}<\dfrac y{14}<\dfrac{10}{14}$, nên $y\in\{7;8;9\}$.
}
\gap
\begin{probbox}{8}{\relax}
Một tấm vải hình chữ nhật có chiều dài $\dfrac74$ m, chiều rộng kém chiều dài $\dfrac5{12}$ m. Tính chu vi và diện tích tấm vải.
\end{probbox}
\probsol{
Chiều rộng: $\dfrac74-\dfrac5{12}=\dfrac{21}{12}-\dfrac5{12}=\dfrac43$ (m).\par
Chu vi: $\left(\dfrac74+\dfrac43\right)\times2=\dfrac{37}{12}\times2=\dfrac{37}6$ (m). \quad Diện tích: $\dfrac74\times\dfrac43=\dfrac73$ (m$^2$).
\dapso{chu vi $\dfrac{37}6$ m; diện tích $\dfrac73$ m$^2$.}}
\gap
\begin{probbox}{9}{\relax}
Một hình chữ nhật có diện tích $\dfrac58$ m$^2$, chiều dài $\dfrac54$ m. Tính chu vi hình chữ nhật đó.
\end{probbox}
\probsol{
Chiều rộng: $\dfrac58:\dfrac54=\dfrac58\times\dfrac45=\dfrac12$ (m). \quad Chu vi: $\left(\dfrac54+\dfrac12\right)\times2=\dfrac72$ (m).
\dapso{$\dfrac72$ m.}}
\gap
\begin{probbox}{10}{\relax}
Từ một sợi dây dài $15$ m, người ta cắt ra $6$ đoạn, mỗi đoạn dài $1\dfrac34$ m. Phần còn lại được cắt thành $3$ đoạn bằng nhau. Hỏi mỗi đoạn sau dài bao nhiêu mét?
\end{probbox}
\probsol{
$6$ đoạn đầu: $1\dfrac34\times6=\dfrac74\times6=\dfrac{21}2$ (m).\par Còn lại: $15-\dfrac{21}2=\dfrac92$ (m).\par
Mỗi đoạn sau: $\dfrac92:3=\dfrac32$ (m).
\dapso{$1\dfrac12$ m.}}
\gap
\begin{probbox}{11}{\relax}
Tính:
\begin{cols2}
$A=\dfrac1{2\times3}+\dfrac1{3\times4}+\dots+\dfrac1{11\times12}$ & $B=\dfrac3{2\times5}+\dfrac3{5\times8}+\dots+\dfrac3{17\times20}$
\end{cols2}
\end{probbox}
\probsol{
\pa{A} $A=\dfrac12-\dfrac13+\dfrac13-\dfrac14+\dots+\dfrac1{11}-\dfrac1{12}=\dfrac12-\dfrac1{12}=\dfrac5{12}$.\par
\pa{B} Mỗi mẫu số là tích của hai số hơn kém nhau $3$, đúng bằng tử số: $B=\dfrac12-\dfrac15+\dfrac15-\dfrac18+\dots+\dfrac1{17}-\dfrac1{20}=\dfrac12-\dfrac1{20}=\dfrac9{20}$.
}
\gap
\begin{probbox}{12}{\relax}
Tính:
\begin{cols2}
$C=\dfrac2{3\times5}+\dfrac2{5\times7}+\dots+\dfrac2{23\times25}$ & $D=\dfrac12+\dfrac14+\dfrac18+\dots+\dfrac1{512}$
\end{cols2}
\end{probbox}
\probsol{
\pa{C} $C=\dfrac13-\dfrac15+\dfrac15-\dfrac17+\dots+\dfrac1{23}-\dfrac1{25}=\dfrac13-\dfrac1{25}=\dfrac{22}{75}$.\par
\pa{D} $2\times D=1+\dfrac12+\dfrac14+\dots+\dfrac1{256}$;\par $2\times D-D=1-\dfrac1{512}$, \ nên $D=\dfrac{511}{512}$.
}
\gap
\begin{probbox}{13}{\relax}
Tìm phân số bằng $\dfrac35$, biết mẫu số hơn tử số $26$ đơn vị.
\end{probbox}
\probsol{
Phân số cần tìm có dạng $\dfrac{3\times k}{5\times k}$.\par Mẫu số trừ tử số bằng $2\times k=26$, nên $k=13$.\par
Phân số cần tìm là $\dfrac{39}{65}$.
\dapso{$\dfrac{39}{65}$.}}
\gap
\begin{probbox}{14}{\relax}
Rút gọn thành phân số tối giản: \ \pa{a} $\dfrac{18}{27}$; \ \pa{b} $\dfrac{56}{70}$; \ \pa{c} $\dfrac{125}{1\,000}$.
\end{probbox}
\probsol{
\pa{a} Chia cả tử số và mẫu số cho $9$: $\dfrac23$.\par \pa{b} Chia cho $14$: $\dfrac45$.\par \pa{c} Chia cho $125$: $\dfrac18$.
\dapso{$\dfrac23$; $\dfrac45$; $\dfrac18$.}}
\gap
\begin{probbox}{15}{\relax}
Quy đồng mẫu số các phân số $\dfrac23,\ \dfrac34,\ \dfrac56$ với mẫu số chung nhỏ nhất.
\end{probbox}
\probsol{
Mẫu số chung nhỏ nhất: $12$.\par $\dfrac23=\dfrac8{12}$;\par $\dfrac34=\dfrac9{12}$;\par $\dfrac56=\dfrac{10}{12}$.
\dapso{$\dfrac8{12};\ \dfrac9{12};\ \dfrac{10}{12}$.}}
\gap
\begin{probbox}{16}{\relax}
Sắp xếp theo thứ tự từ bé đến lớn: $\dfrac35,\ \dfrac23,\ \dfrac47,\ \dfrac7{10}$.
\end{probbox}
\probsol{
Mẫu số chung $210$: $\dfrac{126}{210},\ \dfrac{140}{210},\ \dfrac{120}{210},\ \dfrac{147}{210}$.\par So sánh các tử số: $120<126<140<147$.
\dapso{$\dfrac47<\dfrac35<\dfrac23<\dfrac7{10}$.}}
\gap
\begin{probbox}{17}{\relax}
So sánh $\dfrac5{12}$ và $\dfrac7{13}$ bằng cách dùng $\dfrac12$ làm phân số trung gian.
\end{probbox}
\probsol{
$\dfrac5{12}<\dfrac6{12}=\dfrac12$.\par $\dfrac7{13}>\dfrac{6{,}5}{13}=\dfrac12$.\par Vậy $\dfrac5{12}<\dfrac12<\dfrac7{13}$.
\dapso{$\dfrac5{12}<\dfrac7{13}$.}}
\gap
\begin{probbox}{18}{\relax}
So sánh $\dfrac{2\,023}{2\,024}$ và $\dfrac{2\,024}{2\,025}$.
\end{probbox}
\probsol{
Phần bù đến $1$: $1-\dfrac{2\,023}{2\,024}=\dfrac1{2\,024}$;\par $1-\dfrac{2\,024}{2\,025}=\dfrac1{2\,025}$.\par Vì $\dfrac1{2\,024}>\dfrac1{2\,025}$ nên phân số thứ nhất bé hơn.
\dapso{$\dfrac{2\,023}{2\,024}<\dfrac{2\,024}{2\,025}$.}}
\gap
\begin{probbox}{19}{\relax}
\pa{a} Viết dưới dạng phân số: $3\dfrac25$; \ $5\dfrac38$. \quad \pa{b} Viết dưới dạng hỗn số: $\dfrac{29}6$; \ $\dfrac{47}9$.
\end{probbox}
\probsol{
\pa{a} $3\dfrac25=\dfrac{17}5$;\par $5\dfrac38=\dfrac{43}8$.\par \pa{b} $29=6\times4+5$, nên $\dfrac{29}6=4\dfrac56$;\par $47=9\times5+2$, nên $\dfrac{47}9=5\dfrac29$.
\dapso{$\dfrac{17}5$; $\dfrac{43}8$; $4\dfrac56$; $5\dfrac29$.}}
\gap
\begin{probbox}{20}{\relax}
Tính: \ \pa{a} $3\dfrac14+2\dfrac23$; \ \pa{b} $4\dfrac12-1\dfrac56$.
\end{probbox}
\probsol{
\pa{a} $=\dfrac{13}4+\dfrac83=\dfrac{39}{12}+\dfrac{32}{12}=\dfrac{71}{12}=5\dfrac{11}{12}$.\par \pa{b} $=\dfrac{27}6-\dfrac{11}6=\dfrac{16}6=\dfrac83=2\dfrac23$.
\dapso{$5\dfrac{11}{12}$; $2\dfrac23$.}}
\gap
\begin{probbox}{21}{\relax}
Tính: \ \pa{a} $\dfrac34\times2\dfrac23$; \ \pa{b} $1\dfrac12:\dfrac34$.
\end{probbox}
\probsol{
\pa{a} $=\dfrac34\times\dfrac83=\dfrac{24}{12}=2$.\par \pa{b} $=\dfrac32\times\dfrac43=\dfrac{12}6=2$.
\dapso{$2$; $2$.}}
\gap
\begin{probbox}{22}{\relax}
Tính bằng cách thuận tiện: $\dfrac37\times\dfrac59+\dfrac37\times\dfrac49$.
\end{probbox}
\probsol{
$=\dfrac37\times\left(\dfrac59+\dfrac49\right)$\par $=\dfrac37\times1=\dfrac37$.
\dapso{$\dfrac37$.}}
\gap
\begin{probbox}{23}{\relax}
Tính: $\left(1-\dfrac14\right)\times\left(1+\dfrac13\right)$.
\end{probbox}
\probsol{
$=\dfrac34\times\dfrac43$\par $=1$.
\dapso{$1$.}}
\gap
\begin{probbox}{24}{\relax}
Tìm $x$: \ \pa{a} $x\times\dfrac35=\dfrac9{10}$; \ \pa{b} $x:\dfrac23=\dfrac94$.
\end{probbox}
\probsol{
\pa{a} $x=\dfrac9{10}:\dfrac35=\dfrac9{10}\times\dfrac53=\dfrac32$.\par \pa{b} $x=\dfrac94\times\dfrac23=\dfrac32$.
\dapso{$x=\dfrac32$; $x=\dfrac32$.}}
\gap
\begin{probbox}{25}{\relax}
Tìm $x$: \ \pa{a} $x+\dfrac13=\dfrac56$; \ \pa{b} $\dfrac78-x=\dfrac14$.
\end{probbox}
\probsol{
\pa{a} $x=\dfrac56-\dfrac13=\dfrac56-\dfrac26=\dfrac12$.\par \pa{b} $x=\dfrac78-\dfrac14=\dfrac78-\dfrac28=\dfrac58$.
\dapso{$x=\dfrac12$; $x=\dfrac58$.}}
\gap
\begin{probbox}{26}{\relax}
Một bể đầy nước. Lần thứ nhất người ta dùng $\dfrac25$ bể, lần thứ hai dùng $\dfrac13$ bể. Hỏi trong bể còn lại bao nhiêu phần nước?
\end{probbox}
\probsol{
Đã dùng: $\dfrac25+\dfrac13=\dfrac{11}{15}$ bể.\par Còn lại: $1-\dfrac{11}{15}=\dfrac4{15}$.
\dapso{$\dfrac4{15}$ bể.}}
\gap
\begin{probbox}{27}{\relax}
Một hình chữ nhật có chiều dài $\dfrac56$ m, chiều rộng $\dfrac23$ m. Tính chu vi và diện tích hình chữ nhật đó.
\end{probbox}
\probsol{
Chu vi: $\left(\dfrac56+\dfrac23\right)\times2=\dfrac96\times2=3$ m.\par Diện tích: $\dfrac56\times\dfrac23=\dfrac{10}{18}=\dfrac59$ m$^2$.
\dapso{$3$ m; $\dfrac59$ m$^2$.}}
\gap
\begin{probbox}{28}{\relax}
Tìm phân số bằng $\dfrac35$, biết tổng của tử số và mẫu số bằng $40$.
\end{probbox}
\probsol{
Tử số $3$ phần, mẫu số $5$ phần; $8$ phần bằng $40$.\par Một phần: $40:8=5$.\par Tử số $15$, mẫu số $25$.
\dapso{$\dfrac{15}{25}$.}}
\gap
\begin{probbox}{29}{\relax}
Tìm phân số bằng $\dfrac27$, biết tổng của tử số và mẫu số bằng $54$.
\end{probbox}
\probsol{
$2+7=9$ phần; một phần: $54:9=6$.\par Tử số $12$, mẫu số $42$.
\dapso{$\dfrac{12}{42}$.}}
\gap
\begin{probbox}{30}{\relax}
Một sợi dây dài $7\dfrac12$ m được cắt thành các đoạn, mỗi đoạn dài $\dfrac34$ m. Hỏi cắt được bao nhiêu đoạn?
\end{probbox}
\probsol{
$7\dfrac12=\dfrac{15}2$.\par Số đoạn: $\dfrac{15}2:\dfrac34=\dfrac{15}2\times\dfrac43=10$.
\dapso{$10$ đoạn.}}
\dividerbar{Thử thách -- Nâng cao}
\begin{probbox}{31}{\stars{1}}
Tính:
\begin{cols2}
$E=\dfrac16+\dfrac1{12}+\dfrac1{20}+\dots+\dfrac1{420}$ & $F=\dfrac5{3\times8}+\dfrac5{8\times13}+\dots+\dfrac5{48\times53}$
\end{cols2}
\end{probbox}
\probsol{
\pa{E} $420=20\times21$, nên $E=\dfrac1{2\times3}+\dots+\dfrac1{20\times21}=\dfrac12-\dfrac1{21}=\dfrac{19}{42}$.\par
\pa{F} $F=\dfrac13-\dfrac18+\dfrac18-\dfrac1{13}+\dots+\dfrac1{48}-\dfrac1{53}=\dfrac13-\dfrac1{53}=\dfrac{50}{159}$.
}
\gap
\begin{probbox}{32}{\stars{1}}
Tính $\ G=\dfrac14+\dfrac1{28}+\dfrac1{70}+\dfrac1{130}+\dfrac1{208}$.
\end{probbox}
\probsol{
Nhận xét: $4=1\times4$, $28=4\times7$, $70=7\times10$, $130=10\times13$, $208=13\times16$ (hai thừa số hơn kém nhau $3$).\par Nhân $G$ với $3$:
\[3\times G=\frac3{1\times4}+\frac3{4\times7}+\dots+\frac3{13\times16}=1-\frac1{16}=\frac{15}{16}.\]
Vậy $G=\dfrac{15}{16}:3=\dfrac5{16}$.
}
\gap
\begin{probbox}{33}{\stars{1}}
Tính $\ H=\dfrac23+\dfrac29+\dfrac2{27}+\dfrac2{81}+\dfrac2{243}+\dfrac2{729}$.
\end{probbox}
\probsol{
$3\times H=2+\dfrac23+\dfrac29+\dots+\dfrac2{243}$.\par Trừ đi $H$: $\ 2\times H=2-\dfrac2{729}=\dfrac{1\,456}{729}$, nên $H=\dfrac{728}{729}$.
}
\gap
\begin{probbox}{34}{\stars{1}}
Tính $\ K=\left(1-\dfrac13\right)\times\left(1-\dfrac14\right)\times\left(1-\dfrac15\right)\times\dots\times\left(1-\dfrac1{50}\right)$.
\end{probbox}
\probsol{
\[K=\frac23\times\frac34\times\frac45\times\dots\times\frac{49}{50}=\frac2{50}=\frac1{25}.\]
}
\gap
\begin{probbox}{35}{\stars{2}}
Tính $\ M=\left(1-\dfrac14\right)\times\left(1-\dfrac19\right)\times\left(1-\dfrac1{16}\right)\times\dots\times\left(1-\dfrac1{100}\right)$.
\end{probbox}
\probsol{
Các mẫu số là $2\times2$, $3\times3$, \dots, $10\times10$, và $1-\dfrac1{n\times n}=\dfrac{(n-1)\times(n+1)}{n\times n}$.\par Do đó
\[M=\frac{1\times3}{2\times2}\times\frac{2\times4}{3\times3}\times\frac{3\times5}{4\times4}\times\dots\times\frac{9\times11}{10\times10}.\]
Tách thành hai tích: $\ \dfrac{1\times2\times\dots\times9}{2\times3\times\dots\times10}\times\dfrac{3\times4\times\dots\times11}{2\times3\times\dots\times10}=\dfrac1{10}\times\dfrac{11}2=\dfrac{11}{20}$.
}
\gap
\begin{probbox}{36}{\stars{2}}
Tìm số tự nhiên $y$: $\ \dfrac1{5\times6}+\dfrac1{6\times7}+\dfrac1{7\times8}+\dots+\dfrac1{y\times(y+1)}=\dfrac1{10}$.
\end{probbox}
\probsol{
Vế trái bằng $\dfrac15-\dfrac1{y+1}$.\par Do đó $\dfrac15-\dfrac1{y+1}=\dfrac1{10}$, suy ra $\dfrac1{y+1}=\dfrac15-\dfrac1{10}=\dfrac1{10}$.\par
Vậy $y+1=10$ và $y=9$.
\dapso{$y=9$.}}
\gap
\begin{probbox}{37}{\stars{2}}
Không tính tổng, hãy so sánh $\ S=\dfrac1{11}+\dfrac1{12}+\dfrac1{13}+\dots+\dfrac1{20}\ $ với $\ \dfrac12$.
\end{probbox}
\probsol{
$S$ có $10$ số hạng.\par Mỗi phân số $\dfrac1{11},\dots,\dfrac1{19}$ đều lớn hơn $\dfrac1{20}$, còn số hạng cuối bằng $\dfrac1{20}$.\par
Do đó $S>\dfrac1{20}\times10=\dfrac12$.
\dapso{$S>\dfrac12$.}}
\gap
\begin{probbox}{38}{\stars{1}}
Tính: $\dfrac1{1\times2}+\dfrac1{2\times3}+\dfrac1{3\times4}+\dots+\dfrac1{99\times100}$.
\end{probbox}
\probsol{
Mỗi số hạng $\dfrac1{n\times(n+1)}=\dfrac1n-\dfrac1{n+1}$.\par Tổng $=1-\dfrac12+\dfrac12-\dfrac13+\dots+\dfrac1{99}-\dfrac1{100}$.\par $=1-\dfrac1{100}=\dfrac{99}{100}$.
\dapso{$\dfrac{99}{100}$.}}
\gap
\begin{probbox}{39}{\stars{1}}
Tính: $\dfrac2{1\times3}+\dfrac2{3\times5}+\dfrac2{5\times7}+\dots+\dfrac2{19\times21}$.
\end{probbox}
\probsol{
$\dfrac2{n\times(n+2)}=\dfrac1n-\dfrac1{n+2}$.\par Tổng $=1-\dfrac13+\dfrac13-\dfrac15+\dots+\dfrac1{19}-\dfrac1{21}$.\par $=1-\dfrac1{21}=\dfrac{20}{21}$.
\dapso{$\dfrac{20}{21}$.}}
\gap
\begin{probbox}{40}{\stars{1}}
Tính: $\dfrac12+\dfrac14+\dfrac18+\dots+\dfrac1{1\,024}$.
\end{probbox}
\probsol{
Viết mỗi số hạng thành một hiệu: $\dfrac12=1-\dfrac12$;\par $\dfrac14=\dfrac12-\dfrac14$; \ldots;\par $\dfrac1{1\,024}=\dfrac1{512}-\dfrac1{1\,024}$.\par Cộng lại: $1-\dfrac1{1\,024}=\dfrac{1\,023}{1\,024}$.
\dapso{$\dfrac{1\,023}{1\,024}$.}}
\gap
\begin{probbox}{41}{\stars{1}}
Tính: $\left(1+\dfrac12\right)\times\left(1+\dfrac13\right)\times\left(1+\dfrac14\right)\times\dots\times\left(1+\dfrac1{99}\right)$.
\end{probbox}
\probsol{
$=\dfrac32\times\dfrac43\times\dfrac54\times\dots\times\dfrac{100}{99}$.\par Mỗi tử số rút gọn với mẫu số của phân số liền sau, còn lại $\dfrac{100}2=50$.
\dapso{$50$.}}
\gap
\begin{probbox}{42}{\stars{1}}
So sánh $\dfrac{13}{27}$ và $\dfrac{27}{55}$ mà không quy đồng mẫu số.
\end{probbox}
\probsol{
Cả hai phân số đều bé hơn $\dfrac12$.\par Xét xem mỗi phân số còn kém bao nhiêu:\par $\dfrac12-\dfrac{13}{27}=\dfrac{27-26}{54}=\dfrac1{54}$;\par $\dfrac12-\dfrac{27}{55}=\dfrac{55-54}{110}=\dfrac1{110}$.\par Vì $\dfrac1{54}>\dfrac1{110}$ nên $\dfrac{13}{27}$ kém $\dfrac12$ nhiều hơn, do đó nó bé hơn.
\dapso{$\dfrac{13}{27}<\dfrac{27}{55}$.}}
\gap
\begin{probbox}{43}{\stars{1}}
Viết ba phân số khác nhau lớn hơn $\dfrac13$ và bé hơn $\dfrac12$.
\end{probbox}
\probsol{
Với mẫu số $24$: $\dfrac13=\dfrac8{24}$ và $\dfrac12=\dfrac{12}{24}$.\par Ba phân số: $\dfrac9{24},\ \dfrac{10}{24},\ \dfrac{11}{24}$.
\dapso{Chẳng hạn: $\dfrac9{24};\ \dfrac{10}{24};\ \dfrac{11}{24}$.}}
\gap
\begin{probbox}{44}{\stars{1}}
Tìm tất cả các số tự nhiên $x$ sao cho $\dfrac13<\dfrac x{12}<\dfrac34$.
\end{probbox}
\probsol{
$\dfrac13=\dfrac4{12}$;\par $\dfrac34=\dfrac9{12}$.\par Do đó $4<x<9$, suy ra $x=5,\ 6,\ 7,\ 8$.
\dapso{$x=5;\ 6;\ 7;\ 8$.}}
\gap
\begin{probbox}{45}{\stars{1}}
Phải cộng cùng một số tự nhiên nào vào cả tử số và mẫu số của $\dfrac5{11}$ để được phân số bằng $\dfrac35$?
\end{probbox}
\probsol{
Cộng cùng một số vào tử số và mẫu số thì hiệu $11-5=6$ không đổi.\par Phân số mới bằng $\dfrac35$: hiệu $2$ phần bằng $6$, nên một phần là $3$.\par Tử số mới là $9$; số đã cộng thêm là $9-5=4$.
\dapso{$4$.}}
\gap
\begin{probbox}{46}{\stars{1}}
Phải bớt ở tử số của $\dfrac{23}{37}$ và thêm vào mẫu số cùng một số tự nhiên nào để được phân số bằng $\dfrac13$?
\end{probbox}
\probsol{
Tổng của tử số và mẫu số không đổi: $23+37=60$.\par Phân số mới bằng $\dfrac13$: tử số $1$ phần trong tổng số $4$ phần; tử số mới là $60:4=15$.\par Số cần tìm là $23-15=8$.
\dapso{$8$.}}
\gap
\begin{probbox}{47}{\stars{1}}
Tính: $\dfrac13+\dfrac1{15}+\dfrac1{35}+\dfrac1{63}+\dfrac1{99}$.
\end{probbox}
\probsol{
Viết các mẫu số thành tích: $\dfrac1{1\times3}+\dfrac1{3\times5}+\dfrac1{5\times7}+\dfrac1{7\times9}+\dfrac1{9\times11}$.\par Gấp đôi tổng này được $1-\dfrac1{11}=\dfrac{10}{11}$.\par Vậy tổng bằng $\dfrac{10}{11}:2=\dfrac5{11}$.
\dapso{$\dfrac5{11}$.}}
\gap
\begin{probbox}{48}{\stars{1}}
Viết $\dfrac7{12}$ thành tổng của hai phân số khác nhau có tử số là $1$.
\end{probbox}
\probsol{
Tách $7$ thành tổng hai ước của $12$: $7=3+4$.\par $\dfrac7{12}=\dfrac3{12}+\dfrac4{12}=\dfrac14+\dfrac13$.
\dapso{$\dfrac7{12}=\dfrac13+\dfrac14$ (hoặc $\dfrac12+\dfrac1{12}$).}}
\gap
\begin{probbox}{49}{\stars{2}}
Tính: $1+\dfrac13+\dfrac16+\dfrac1{10}+\dfrac1{15}+\dots+\dfrac1{45}$ (các mẫu số là $1,\ 3,\ 6,\ 10,\ 15,\ \dots$).
\end{probbox}
\probsol{
Mỗi số hạng có dạng $\dfrac2{n\times(n+1)}$: $1=\dfrac2{1\times2}$, $\dfrac13=\dfrac2{2\times3}$, \ldots, $\dfrac1{45}=\dfrac2{9\times10}$.\par Tổng $=2\times\left(1-\dfrac1{10}\right)=\dfrac95$.
\dapso{$\dfrac95=1\dfrac45$.}}
\gap
\begin{probbox}{50}{\stars{2}}
Tìm phân số có mẫu số là $13$, biết rằng nếu cộng thêm $6$ vào tử số (giữ nguyên mẫu số) thì được phân số mới gấp $3$ lần phân số đã cho.
\end{probbox}
\probsol{
Mẫu số không đổi nên tử số mới gấp $3$ lần tử số cũ.\par Do đó $6$ gấp $3-1=2$ lần tử số cũ, tử số cũ là $6:2=3$.\par Phân số cần tìm là $\dfrac3{13}$.
\dapso{$\dfrac3{13}$.}}
\clearpage
\titlehead
